\chapter{State of the Art}
\label{ch:state_of_the_art}

\section{Streaming and Event-Time Foundations}

The architecture is grounded in event-time stream processing and replayable logs, not ingestion-time shortcuts. Core references are the Dataflow model and Flink/Kafka stream infrastructure \citep{akidau2015dataflow,carbone2015flink,kreps2011kafka}.

\begin{itemize}
    \item Event-time ordering and watermarks for causal processing.
    \item Stateful keyed operators for race-context accumulation.
    \item Replayability for reproducible experimental regeneration.
\end{itemize}

% TODO: expand prose and include one short comparative paragraph on event-time vs processing-time.

\section{Race Strategy Decision Support Literature}

Race strategy support combines tactical heuristics, simulation/optimization, and recent data-driven approaches \citep{heilmeier2020race,carrasco2023optimization,quiroga2024optimization,piccinotti2020openloop,widanage2019indy500}.

\begin{itemize}
    \item Optimization-oriented work formalizes undercut/overcut and stop-window trade-offs.
    \item Data-driven work often improves prediction, but reproducibility and evaluation-contract alignment are frequently under-specified.
\end{itemize}

% TODO: classify cited papers by objective (timing vs utility vs simulation) in final prose.

\section{Batch Learning for Tabular Racing Context}

Batch XGBoost is a strong baseline for heterogeneous tabular features under skewed labels \citep{chen2016xgboost}. In this thesis, its relevance is methodological:
\begin{itemize}
    \item robust non-linear feature interactions,
    \item compatibility with structured cross-validation,
    \item direct TreeSHAP interpretability path.
\end{itemize}

Related motorsport ML references should be used as context, not as replacements for method-grounding citations \citep{sasikumar2025pitstop,hettmann2024datatopodium,dinh2025gradient}.

\section{Online Learning and Prequential Evaluation}

MOA provides the stream-mining framework and prequential evaluation mechanics used here \citep{bifet2010moa}. Drift literature frames why online adaptation matters \citep{gama2014drift,bifet2007adwin,gomes2017arf,arora2024driftreview}.

\begin{itemize}
    \item Prequential loop: test on each instance, then train on that instance.
    \item Stream metrics are sensitive to ordering, support, and drift regime.
\end{itemize}

% TODO: if ARF/ADWIN internals are not central in final code path, keep them as contextual citations only.

\section{Imbalanced Evaluation and Threshold Policy}

Under rare positives, PR/AP diagnostics are more informative than ROC-focused summaries \citep{saito2015prroc,davis2006relationship}. Cost-sensitive threshold policy is an explicit design choice, not a by-product \citep{elkan2001costsensitive}.

For this thesis:
\begin{itemize}
    \item Kappa is used as chance-corrected diagnostic agreement \citep{cohen1960kappa}.
    \item G-Mean is used as imbalance-aware class-balance diagnostic \citep{he2009imbalanced}.
\end{itemize}

% TODO: add one compact metric-purpose table in final writing.

\section{Leakage and Structured Validation}

Temporal and grouped split discipline is essential when rows share race context and future outcomes \citep{roberts2017cvstructured,brookshire2024leakage}.

\begin{itemize}
    \item Structured OOF validation is required for fair downstream frontier selection.
    \item Feature filtering must remove target/truth/matched-event shortcuts.
\end{itemize}

\section{Explainability and Probability-Quality Caveats}

Explainability references:
\begin{itemize}
    \item SHAP additive attribution \citep{lundberg2017shap}.
    \item TreeSHAP for tree ensembles (Batch XGBoost) \citep{lundberg2018treeshap}.
\end{itemize}

Calibration/probability-quality references:
\begin{itemize}
    \item probability quality vs ranking quality distinction \citep{niculescu2005goodprob},
    \item optional calibration caveat context \citep{guo2017calibration}.
\end{itemize}

% TODO: keep calibration subsection short unless explicit calibration diagnostics are central to final claims.

\section{Research Gap Addressed by This Thesis}

\begin{itemize}
    \item End-to-end replay-to-evaluation reproducibility.
    \item Shared truth-universe evaluation across deterministic, batch, and streaming paradigms.
    \item Explicit dual-ledger accounting (row decisions vs event coverage).
    \item Corrected strict operational \texttt{pit\_success\_h2} contract and threshold-policy interpretation.
\end{itemize}

\paragraph{Writing note.}
Avoid over-citation. Use core method citations for claims that depend on them, and keep recent motorsport references as contextual support unless directly reused in implementation.
